% figure (2026-11-17): the e- mu- -> e- mu- diagram, time upward, after PHY 513 Lecture 21 slide 4; generated by qft-C/tools/l21/figs/make_figs.sh
\tikzset{
  fdl/.style={line width=0.8pt, black},
  fdx/.style={circle, fill=black, inner sep=0pt, minimum size=3.6pt},
  fdv/.style={circle, fill=figB, draw=figB, inner sep=0pt, minimum size=5pt},
  fdlab/.style={font=\scriptsize, inner sep=2.6pt},
  fdmom/.style={->, >=Stealth, figO, line width=0.6pt},
}
\tikzset{
  fda/.style={fdl, postaction={decorate}, decoration={markings, mark=at position 0.55 with {\arrow{Stealth[length=5pt]}}}},
  fdS/.style={font=\small},
}
\tikzset{
  fline/.style={line width=0.8pt, black, postaction={decorate}, decoration={markings, mark=at position 0.56 with {\arrow{Stealth[length=5.5pt, width=4.5pt]}}}},
  sline/.style={line width=0.8pt, black, dash pattern=on 3.5pt off 2.2pt},
  blob/.style={draw, line width=0.8pt, fill=figB!12},
}
\tikzset{
  pline/.style={line width=0.8pt, black, decorate, decoration={snake, amplitude=1.6pt, segment length=5.2pt}},
}
\begin{tikzpicture}[font=\small]
  \coordinate (v1) at (0,0); \coordinate (v2) at (3.0,0);
  \draw[fline] (-1.7,-1.7)--(v1);
  \draw[fline] (v1)--(-1.7,1.7);
  \draw[fline] (4.7,-1.7)--(v2);
  \draw[fline] (v2)--(4.7,1.7);
  \draw[pline] (v1)--(v2); \node[fdv] at (v1) {}; \node[fdv] at (v2) {};
  % momentum arrows (all along time)
  \draw[fdmom] (-1.02,-1.42)--(-0.52,-0.92); \draw[fdmom] (-0.92,0.52)--(-1.42,1.02);
  \draw[fdmom] (4.02,-1.42)--(3.52,-0.92); \draw[fdmom] (3.92,0.52)--(4.42,1.02);
  \node[fdlab, left] at (-1.7,-1.7) {$e^-\ (p_1, s)$};
  \node[fdlab, left] at (-1.7,1.7) {$e^-\ (p_1', s')$};
  \node[fdlab, right] at (4.7,-1.7) {$\mu^-\ (p_2, r)$};
  \node[fdlab, right] at (4.7,1.7) {$\mu^-\ (p_2', r')$};
  \draw[fdmom] (1.1,0.36)--(1.9,0.36); \node[fdlab, above] at (1.5,0.36) {$q = p_1 - p_1'$};
  \node[fdlab, below] at (1.5,-0.16) {$\gamma^*$, off shell};
  \node[fdlab, below] at (0.05,-0.16) {$\mu$}; \node[fdlab, below] at (2.95,-0.16) {$\nu$};
  \node[fdlab] at (-2.0,-0.75) {$u^s(p_1)$}; \node[fdlab] at (-2.0,0.75) {$\bar u^{s'}(p_1')$};
  \node[fdlab] at (5.0,-0.75) {$u^r(p_2)$}; \node[fdlab] at (5.0,0.75) {$\bar u^{r'}(p_2')$};
  \draw[->, gray] (6.9,-1.2)--(6.9,1.2); \node[font=\scriptsize, gray, right] at (6.9,0) {time};
  \node[fdS, align=center] at (1.5,-2.75) {$i\mathcal M = \underbrace{\bigl[\bar u^{s'}(p_1')(-ie\gamma^\mu)u^s(p_1)\bigr]}_{\text{electron line}}\ \underbrace{\dfrac{-ig_{\mu\nu}}{q^2}}_{\text{propagator}}\ \underbrace{\bigl[\bar u^{r'}(p_2')(-ie\gamma^\nu)u^r(p_2)\bigr]}_{\text{muon line}}$};
\end{tikzpicture}
